%% Abstract + keywords (SPIE style).
%% Rev. 2 --- adapted from the pre-study draft to what the results actually show:
%% the RL policy does NOT dominate; it ties two very different architectures in the
%% realistic deployment condition. Every number is traceable to
%% papers/paper1-benchmark/README.md and the reports/ notebook.

\begin{abstract}
Quadrotors leaving controlled conditions must track trajectories under sustained and
unsteady wind, under changes in mass, and on state estimates far poorer than motion
capture. Which controller to choose for that regime is an open practical question:
candidates are published on different vehicles, references, metrics and sensing stacks,
so their relative merits, and especially their failure modes, cannot be compared across
papers. We benchmark six trajectory-tracking controllers, one per control paradigm, on a
common testbed --- one simulated airframe, one reference family
(figure-eight Lissajous at three cycle times) and a single command interface: a PID with
acceleration feedforward, active disturbance rejection control with an extended state
observer, sampling-based MPPI with $\mathcal{L}_1$ adaptation, an offset-free nonlinear
MPC, a hybrid stack pairing the observer with a pretrained learned adaptive rate
controller, and a reinforcement-learning policy trained with a privileged critic and
sensor-noise domain randomization. The testbed is a brushless Crazyflie~2.1 ($43.4$~g,
thrust-to-weight $1.88$), chosen because a tight authority margin and modest onboard
sensing make the trade-offs under study large enough to measure cleanly; the absolute
errors are specific to it, while the mechanisms they expose are not. All six are
evaluated under constant wind, gusts,
in-ground effect, a $23\,\%$ payload step, and a Lighthouse measurement model grounded
in published characterization data, using ten evaluation seeds and three training seeds
so that every claim carries an error bar. No controller dominates: the per-condition
champions differ, and each wins through a different disturbance-estimation mechanism.
Under the realistic deployment condition --- degraded onboard sensing with wind --- four
of them become statistically indistinguishable, spanning $0.057$ to $0.065$~m: an
offset-free MPC, an extended state observer over a learned adaptive rate loop, that same
observer alone, and the learned policy. The obvious reading is that the shared state
estimate is what binds. We tested that reading by scaling every sensor error term
together, and it does not survive: a fourfold better sensor improves the four stacks by
only $3$ to $16\,\%$ and leaves them tied, while a twofold worse one separates them and
triples their spread. Measurement quality bounds this regime from above without setting
its floor, so sensor upgrades protect against the tail rather than buying accuracy. We
distil the study into seven transferable lessons. Three of them record occasions on
which our own earlier reading of the data was wrong and the protocol caught it: a launch
transient misdiagnosed as noise fragility, a single evaluation seed that proved to be the
minimum of ten draws, and a three-seed variance figure that overstated both its mean and
its spread. Finally, we describe the hardware
validation protocol for the recommended controller --- a programmable fan-array wind
tunnel with motion-capture ground truth --- designed so that the physical disturbance
conditions correspond one-to-one with the simulated ones.
% TODO (after the wind-tunnel campaign): replace the last sentence with the measured
% sim-to-real outcome, e.g. "...; in flight, the recommended stack tracked to XX m
% under a YY m/s programmed gust field, an increase of ZZx over its simulated value."
\end{abstract}

\keywords{quadrotor control, trajectory tracking, benchmarking, disturbance rejection,
active disturbance rejection control, model predictive control, reinforcement learning,
domain randomization, wind tunnel, motion capture, unmanned aerial vehicles}
